\begin{abstract}
We prove that a discrete group is amenable if and only if every
uniformly bounded representation on a complex Hilbert space is similar
to a unitary representation. This resolves Dixmier's unitarizability
problem affirmatively for discrete groups.
\end{abstract}

\section{Introduction}\label{sec:introduction}

A representation \(\pi:G\to\mathrm{GL}(H)\) of a discrete group on a
complex Hilbert space is \emph{uniformly bounded} if
\[
 |\pi|:=\sup_{g\in G}\norm{\pi(g)}<\infty.
\]
Here \(\mathrm{GL}(H)\) consists of the bounded invertible operators.
The representation is \emph{unitarizable} if some bounded invertible
\(S\) makes every \(S\pi(g)S^{-1}\) unitary. Equivalently, the Hilbert
space admits an equivalent inner product invariant under the
representation. A group is unitarizable when every uniformly bounded
representation on every complex Hilbert space is unitarizable. The
similarity may depend on the representation; its condition number
\(\norm{S}\norm{S^{-1}}\) is not assumed to have a common bound.

A discrete group \(G\) is \emph{amenable} if there is a positive linear
functional \(m:\ell^\infty(G)\to\C\) with \(m(1)=1\) and
\(m(L_gf)=m(f)\), where \((L_gf)(x)=f(g^{-1}x)\).
The one-operator case goes back to Sz.-Nagy: an invertible operator whose
integer powers are uniformly bounded is similar to a unitary operator
\cite[Theorem~I]{SzNagy1947}. His generalized-limit construction of an
invariant inner product preceded the invariant-mean argument of Day and
Dixmier \cite{Day1950,Dixmier1950}. Thus amenability implies
unitarizability. Dixmier explicitly asked whether this property
characterizes groups admitting an invariant mean
\cite[\S5, p.~221]{Dixmier1950}. The question asks whether changing the
Hilbert norm separately for every uniformly bounded representation
already forces the existence of such a mean.

\begin{theorem}\label{thm:main}
A discrete group is amenable if and only if it is unitarizable.
For every discrete nonamenable group \(G\) and every \(\varepsilon>0\),
there are a complex Hilbert space \(\mathcal K\) and a representation
\(\pi:G\to\mathrm{GL}(\mathcal K)\) with \(|\pi|\le1+\varepsilon\)
that is not similar to a unitary representation. If \(G\) is countable,
\(\mathcal K\) may be chosen separable.
\end{theorem}

We work in ZFC. All Hilbert spaces below are complex, and
\(\N=\{1,2,\ldots\}\). The main point is the absence of a similarity for one
representation, without any a priori control on possible similarities.
For an uncountable group the induced witness may be nonseparable, as is
allowed by the definition of unitarizability used here.

\subsection*{Prior work and the remaining obstruction}

Early counterexamples showed that uniform boundedness alone does not
force unitarizability. Ehrenpreis and Mautner constructed examples for
the locally compact group \(\mathrm{SL}_2(\mathbb R)\) \cite{EhrenpreisMautner1955}; explicit
families for free groups include those of Pytlik and Szwarc
\cite{PytlikSzwarc1986}. Nonunitarizability passes from a subgroup to the
ambient group by induction. The free-group examples therefore cover groups
containing nonabelian free subgroups, but leave nonamenable groups
without such subgroups to be understood.

Quantitative unitarizability gives one route to amenability. Pisier
proved that a discrete group \(G\) is amenable exactly when there are
constants \(C\) and \(\alpha<3\), depending only on \(G\), such that
every uniformly bounded representation admits a unitarizer with
condition number at most \(C|\pi|^\alpha\)
\cite[Theorem~0.1]{Pisier1998}. The estimate is uniform over all Hilbert
spaces and representations of the fixed group. Unitarizability itself
implies a polynomial bound of this form for some finite exponent
\cite[Theorem~4.1]{Pisier2004}; the additional content of the amenability
criterion is the threshold \(\alpha<3\).

Random forests gave a route beyond free subgroups. Epstein and Monod
used invariant random forests to relate unitarizability to
\(\ell^2\)-Betti numbers and cost. Their estimates
imply nonunitarizability for residually finite groups with positive
first \(\ell^2\)-Betti number, and for finitely generated residually
finite groups of cost greater than one
\cite[Theorems~1.1--1.3]{EpsteinMonod2009}. Combining their criterion
with group constructions, Osin obtained finitely generated torsion and
torsion-free nonunitarizable groups containing no nonabelian free
subgroup \cite[Theorems~1.2--1.3]{Osin2009}.

A different approach uses measurable actions. Gaboriau and Lyons showed
that a Bernoulli action of a countable nonamenable group contains,
within its orbits, those of an essentially free probability-preserving
action of the free group \(F_2\) \cite[Theorem~1]{GaboriauLyons2009}.
Monod and Ozawa used this to show that, for an
infinite abelian group \(A\), the restricted wreath product
\(A\wr G=(\bigoplus_G A)\rtimes G\) is unitarizable exactly when
\(G\) is amenable \cite[Theorem~1]{MonodOzawa2010}. This places the
obstruction on an extension of \(G\). Alpeev subsequently allowed every
nontrivial countable lamp group \(A\) when \(G\) is countable and
nonamenable, so the conclusion also holds for finite lamps
\cite[Theorem~A]{Alpeev2023}. The distinction remains between a witness
for an extension and one for \(G\) itself.

Approximation properties give a complementary route. Vergara proved
that a discrete unitarizable group satisfying the
\(M_d\)-approximation property for every \(d\ge2\) is amenable.
These conditions hold for countable groups acting properly on
finite-dimensional CAT(0) cube complexes
\cite[Theorems~1.2--1.3]{Vergara2023}. His later extension theorem
preserves \(M_d\)-AP when both a normal subgroup and the quotient
possess it \cite[Theorem~1.3]{Vergara2025}.

Littlewood functions give the analytic formulation closest to the
present proof. Their scalar convolution kernels split into a part with
uniformly absolutely summable rows and a part with uniformly absolutely
summable columns. Epstein and Monod, following Bo\.{z}ejko and Fendler,
use this splitting to obtain bounded commutators and triangular
representations \cite[Propositions~2.3--2.4]{EpsteinMonod2009}.
Gerasimova, Gruber, Monod and Thom relate the corresponding norms to
Cheeger constants and study the resulting restrictions on unitarizable
groups \cite[\S2.B]{GGMT2018}. The passage from bounded non-inner
operator cocycles to nonunitarizable representations is classical
\cite[\S2]{MonodOzawa2010},
\cite[Remark~4 and Proposition~5(ii)]{Pisier2011}.
The present construction retains that mechanism while allowing
finite-dimensional matrix coefficients. Positive rank-one weights
preserve the row and column bounds while increasing the kernel's
Hilbert--Schmidt energy relative to the fibre dimension. This gives a
direct obstruction with no prescribed bound on possible unitarizers.

The passage from countable groups to arbitrary discrete groups uses
standard facts. F\o lner's finite-set criterion implies that every
nonamenable discrete group has a finitely generated nonamenable subgroup
\cite[Main Theorem, p.~245]{Folner1955}.
Induction carries a nonunitarizable representation of that subgroup
to the full group with the same uniform bound
\cite[Proposition~0.5]{Pisier2004}; Pisier's Corollaries~0.10--0.11 also
show that unitarizability of discrete groups is countably determined.

\subsection*{The operator obstruction and the proof}

Let \(U\) be a unitary representation of \(G\) on a Hilbert space
\(\mathcal H\), and write \(\mathcal B(\mathcal H)\) for its bounded
operators. A bounded operator cocycle is a map
\(D:G\to\mathcal B(\mathcal H)\) with
\(\sup_g\norm{D(g)}<\infty\) and
\[
 D(gh)=D(g)+U(g)D(h)U(g)^{-1}.
\]
It is \emph{inner} if one bounded operator \(B\) satisfies
\(D(g)=B-U(g)BU(g)^{-1}\) for all \(g\).
The cocycle gives the uniformly bounded triangular representation
\[
 \pi(g)=\begin{pmatrix}U(g)&D(g)U(g)\\0&U(g)\end{pmatrix}.
\]
A similarity making this representation unitary would produce an
invariant complement to \(\mathcal H\oplus0\). Writing that complement
as the graph of a bounded operator makes \(D\) inner.
Thus it suffices to construct a bounded cocycle with no bounded
implementer.

The criterion in \Cref{sec:obstruction} reduces this task to operators
\(T_j\) on \(\ell^2(G;\C^{k_j})\). Their conjugation differences under
left translation must be bounded uniformly in both \(j\) and \(g\),
while their distance from the operators commuting with left translation
must tend to infinity. A translation-invariant comparison operator
\(A_j\), with blocks \(A_j[x,y]=W_j(x^{-1}y)\), detects that distance:
if the block row at \(e\) of \(T_j\) and the block column at \(e\) of \(A_j-T_j\) have
norm at most \(K\), then
\[
 \dist(T_j,U_j(G)')\ge
 \frac12\left(\frac{\sum_{s\in G}\norm{W_j(s)}_{\HS}^2}{k_j}\right)^{1/2}-K.
\]
Here \(U_j\) is left translation, \(U_j(G)'\) is its commutant, and
\(\norm{W}_{\HS}^2=\operatorname{Tr}(W^*W)\) uses the unnormalized
trace. This estimate compares one finite-dimensional row with one
finite-dimensional column; it imposes no bound on \(\norm{T_j}\).

Nonamenability supplies the operators in two steps. In
\Cref{sec:sparse}, a contracting finite average gives lists
\(s_1,\ldots,s_n\) of group elements, with multiplicity. For \(x\in G\)
and \(1\le i\le n\), the labelled edge \((x,i)\) joins row \(x\) to
column \(xs_i\). These edges can be
assigned to their endpoints so that each endpoint receives at most
\(r\) edges, where \((r/n)\log n\to0\). Compare an assignment with
any left translate. At a row, the edges whose assignment changes lie
in the two sets assigned to that row; at a column, they lie in the
two sets assigned to that column. Thus at most \(2r\) labels change
at either endpoint. In
\Cref{sec:weights}, we choose \(n\) unit vectors in \(\C^k\), with
\(k\) of order \(r\log n\), whose synthesis operators have bounded
norm on every subfamily of at most \(2r\) vectors. This is the
concentration, finite-net and union-bound method familiar from random
matrix estimates; compare \cite[Section~5]{BDDW2008}.
Only its upper norm bound is needed here.

The rank-one projections onto these vectors weight the convolution
edges, and \(T_j\) retains the row-assigned edges. Factoring through
the scalar Hilbert space indexed by the edges
bounds the relevant rows, columns and all conjugation differences.
Keeping the edge labels is essential: distinct products may coincide
in \(G\), yet correspond to distinct coordinates in the factorization.
Positivity ensures that such coincidences do not reduce the kernel
energy, which is at least \(n\). Since \(k/n\to0\), the distance
estimate diverges. Finally, \Cref{sec:conclusion} rescales these operators
and applies the
obstruction to one direct sum of their cocycles, giving any prescribed
uniform bound \(1+\varepsilon\). Induction passes the witness from a
finitely generated nonamenable subgroup to an arbitrary discrete group.
The Day--Dixmier averaging argument completes the equivalence.
\Cref{sec:cstar-consequence} then uses the separate companion similarity
theorem to show that the resulting \(\ell^1(G)\)-representations have no
bounded algebra-homomorphism extension to either group \(C^*\)-algebra.
